\section{References}

% Long archival URLs cannot always justify cleanly; the reference lists are
% therefore set ragged right within this section only.
\begingroup
\raggedright
% One archival URL below is a single hyphenated segment longer than the
% measure; within this section only, URLs may also break after hyphens.
\expandafter\def\expandafter\UrlBreaks\expandafter{\UrlBreaks\do\-}

All web sources were accessed and checked on 2026-07-25. Citations in the
text use conventional titles and standard internal divisions.

\subsection*{Primary sources: Thomas Aquinas}

\begin{itemize}
\item \work{Summa theologiae} I.1.1, corpus opening, Latin: \textit{Sancti
Thomae Aquinatis Opera omnia}, tomus quartus (Leonine edition; Rome:
Typographia Polyglotta, 1888), p.~6; wording as centrally verified in this
repository's source library
(\texttt{passage.thomas-aquinas.summa-theologiae.leonine-v4-1888.i-q1-a1-corpus-opening},
verified 2026-07-23) and re-read in the Corpus Thomisticum web
presentation of the Leonine text
(\url{https://www.corpusthomisticum.org/sth1001.html}; database
presentation \copyright~2019 Fundaci\'on Tom\'as de Aquino).
\item \work{Summa theologiae}, English: \textit{The Summa Theologica of
St.\ Thomas Aquinas}, literally translated by Fathers of the English
Dominican Province, 2nd rev.\ ed.\ (1920; public domain), as transcribed
at New Advent (\url{https://www.newadvent.org/summa/1001.htm} for I.1);
the general prologue is paraphrased in this study, not quoted (see
appendix).
\item \work{Summa contra Gentiles} I.2, Latin, at Corpus Thomisticum
(\url{https://www.corpusthomisticum.org/scg1001.html}); brief quotation
with attribution.
\item \work{Summa theologiae} I.2 prologue (plan of the work), English as
quoted in Benedict XVI's General Audience of 23 June 2010 (below).
\item Works cited by title only (\work{Scriptum super Sententiis},
\work{De ente et essentia}, \work{Contra impugnantes}, \work{De
veritate}, \work{Catena aurea}, \work{Contra errores Graecorum},
\work{De unitate intellectus}, \work{De malo}, \work{De virtutibus},
\work{De perfectione spiritualis vitae}, \work{Contra retrahentes}):
identified through the reference literature below; no locus-level claim
rests on an uninspected text.
\end{itemize}

\subsection*{Early biographical witnesses (identified; not inspected)}

\begin{itemize}
\item \textit{Fontes vitae S.\ Thomae Aquinatis, notis historicis et
criticis illustrati}, ed.\ D.~Pr\"ummer and M.-H.~Laurent (Toulouse /
Revue Thomiste, 1911--1937): the canonization processes (Naples 1319;
Fossanova 1321) and the Lives of William of Tocco, Bernard Gui, and
Peter Calo. Scan located at the Internet Archive
(\url{https://archive.org/details/prummerfontesvitaestho};
\url{https://archive.org/details/fontesvitaesthom00pr}); catalog-level
identification only; not page-inspected.
\item William of Tocco, \work{Ystoria sancti Thome de Aquino}
(1318--1323), critical edition: Claire Le Brun-Gouanvic (Toronto: PIMS,
Studies and Texts 127, 1996; ISBN 978-0-88844-127-0)
(\url{https://pims.ca/publication/isbn-978-0-88844-127-0/}); publisher
record only.
\item Kenelm Foster, O.P. (ed., trans.), \textit{The Life of Saint
Thomas Aquinas: Biographical Documents} (London: Longmans, Green;
Baltimore: Helicon, 1959). In copyright; identified through publisher
and review records
(\url{https://openlibrary.org/books/OL6269895M/The_life_of_Saint_Thomas_Aquinas});
not quoted.
\end{itemize}

\subsection*{Reception-era reference (public domain; checked)}

\begin{itemize}
\item D.~J. Kennedy, ``St.\ Thomas Aquinas,'' \textit{The Catholic
Encyclopedia}, vol.~14 (New York, 1912), at New Advent
(\url{https://www.newadvent.org/cathen/14663b.htm}); principal
biographical frame and the quoted renderings of the canonization-era
traditions.
\item D.~J. Kennedy, ``Thomism,'' \textit{The Catholic Encyclopedia},
vol.~14 (New York, 1912), at New Advent
(\url{https://www.newadvent.org/cathen/14698b.htm}); 1277 censures,
Dominican defense, Bourret revocation quotation.
\item ``Corpus Christi, Feast of,'' \textit{The Catholic Encyclopedia},
vol.~4 (New York, 1912), at New Advent
(\url{https://www.newadvent.org/cathen/04390b.htm}); Juliana of Li\`ege,
\work{Transiturus}, 1912-era attribution of the office.
\end{itemize}

\subsection*{Modern scholarship and reference}

\begin{itemize}
\item ``Thomas Aquinas,'' \textit{Stanford Encyclopedia of Philosophy}
(first published 7 December 2022), checked at
\url{https://plato.stanford.edu/entries/aquinas/}; chronology and the
1277 relation.
\item J.~M.~M.~H. Thijssen, ``Condemnation of 1277,'' \textit{Stanford
Encyclopedia of Philosophy} (substantive revision 18 December 2023),
checked at \url{https://plato.stanford.edu/entries/condemnation/};
Tempier's acts, Wippel/Hissette debate, Giles of Rome, Wielockx thesis,
Kilwardby.
\item J.-P. Torrell, O.P., \textit{Saint Thomas Aquinas}, vol.~1:
\textit{The Person and His Work}, trans.\ R.~Royal (Washington: CUA
Press, rev.\ ed.\ 2005; 3rd ed.\ trans.\ Royal and Minerd, ISBN
978-0-8132-3560-8, \url{https://www.cuapress.org/9780813235608/saint-thomas-aquinas/}).
Publisher-record level; not independently inspected.
\item M.~Grabmann, \textit{Thomas Aquinas: His Personality and Thought},
trans.\ V.~Michel (New York: Longmans, Green, 1928; repr.\ 1963).
Catalog-record level; not independently inspected.
\item M.~R\"as\"anen, \textit{Thomas Aquinas's Relics as Focus for
Conflict and Cult in the Late Middle Ages: The Restless Corpse}
(Amsterdam: Amsterdam University Press, 2017; ISBN 978-90-8964-873-0,
\url{https://www.aup.nl/en/book/9789089648730/thomas-aquinas-s-relics-as-focus-for-conflict-and-cult-in-the-late-middle-ages}).
Publisher-record level; not independently inspected.
\item C.~Lambot, ``L'office de la F\^ete-Dieu. Aper\c{c}us nouveaux sur
ses origines,'' \textit{Revue b\'en\'edictine} (1942), and P.-M. Gy,
``L'office du Corpus Christi et la th\'eologie des accidents
eucharistiques,'' \textit{Revue des sciences philosophiques et
th\'eologiques} (1982). Bibliographic identification only; frame of the
office-attribution debate.
\item Commissio Leonina, site checked at
\url{https://www.commissio-leonina.org/}; the commission's continuing
critical-edition charge.
\item ``The Relics of the Angelic Doctor,'' Prime Matters (University of
Mary), checked at
\url{https://primematters.com/foundations/enrichment/relics-angelic-doctor};
reported level for the Revolution-era transfer, 1974 return, and Naples
relic claims.
\end{itemize}

\subsection*{Ecclesial reception (checked at the Holy See's site except
as noted)}

\begin{itemize}
\item Leo XIII, encyclical \textit{Aeterni Patris} (4 August 1879),
English text
(\url{https://www.vatican.va/content/leo-xiii/en/encyclicals/documents/hf_l-xiii_enc_04081879_aeterni-patris.html});
quoted briefly. Distinct from Pius IX's 1868 bull \textit{Aeterni
Patris} convoking the First Vatican Council, which is not used here.
\item Sacred Congregation of Studies, decree approving twenty-four
Thomistic theses (27 July 1914), English at the Franciscan Archive
(\url{https://franciscan-archive.org/thomas/24theses.html}); preamble
reporting Pius X's motu proprio \textit{Doctoris Angelici} (29 June
1914), which is otherwise cited at reported level only.
\item Pius XI, encyclical \textit{Studiorum Ducem} (29 June 1923), Latin
text
(\url{https://www.vatican.va/content/pius-xi/la/encyclicals/documents/hf_p-xi_enc_19230629_studiorum-ducem.html});
``Doctor Communis'' and the 1917 canon 1366~\S2 quotation.
\item Second Vatican Council, decree \textit{Optatam totius} (28 October
1965) 16, English text
(\url{https://www.vatican.va/archive/hist_councils/ii_vatican_council/documents/vat-ii_decree_19651028_optatam-totius_en.html}).
\item Paul VI, apostolic letter \textit{Lumen Ecclesiae} (20 November
1974), Italian text
(\url{https://www.vatican.va/content/paul-vi/it/apost_letters/documents/hf_p-vi_apl_19741205_lumen-ecclesiae.html});
paraphrased.
\item \textit{Code of Canon Law} (1983), canon 252~\S3, English text
(\url{https://www.vatican.va/archive/cod-iuris-canonici/eng/documents/cic_lib2-cann208-329_en.html}).
\item \textit{Catechism of the Catholic Church} 460, 1381, 2763, English
text, Vatican website archive
(\url{https://www.vatican.va/archive/ENG0015/__P1J.HTM};
\url{https://www.vatican.va/archive/ENG0015/__P41.HTM};
\url{https://www.vatican.va/archive/ENG0015/__P9X.HTM}).
\item John Paul II, encyclical \textit{Fides et ratio} (14 September
1998) 43--44, English text
(\url{https://www.vatican.va/content/john-paul-ii/en/encyclicals/documents/hf_jp-ii_enc_14091998_fides-et-ratio.html}).
\item Benedict XVI, General Audiences of 2, 16, and 23 June 2010 (three
catecheses on Saint Thomas Aquinas), English texts
(\url{https://www.vatican.va/content/benedict-xvi/en/audiences/2010/documents/hf_ben-xvi_aud_20100602.html};
\ldots\texttt{20100616.html}; \ldots\texttt{20100623.html});
paraphrased, with the Summa plan quotation as cited there.
\item Leo XIV's declaration of John Henry Newman as Doctor of the Church
and co-patron, with Thomas Aquinas, of the Church's educational mission
(1 November 2025): reported from G.~O'Connell, America Magazine, 1
November 2025
(\url{https://www.americamagazine.org/vatican-dispatch/2025/11/01/pope-leo-declares-st-john-henry-newman-a-doctor-of-the-church-and-co-patron-of-catholic-education/});
the formal act's text was not checked.
\end{itemize}

\endgroup
